\clearpage
\section{Critic: what the readings leave unresolved}
\label{sec:critic}

The five readings are useful precisely because their weaknesses are
different. Rather than criticising each paper at length, I reduce the issues
to four questions that matter for implementation.

\subsection{Where did the threshold come from?}

Grubbs' critical value is theoretically defined once a significance level is
chosen, but it assumes a particular statistical setting. The modified
z-score threshold 3.5 is an empirical screening recommendation for batches
\citep{iglewicz1993}. Operational range, step and persistence limits are
often engineering judgements \citep{shafer2000}. In practice these numbers
can survive longer than the reasons behind them.

My experiment makes this concrete. A threshold can reduce false alarms by
sacrificing detection, but there is no single number that fixes a mismatch
between the method and the data. A deployment should therefore record the
source of every threshold, the sampling rate and window for which it was
chosen, and the measured false-alarm/missed-fault trade-off.

\subsection{What does the reviewer experience?}

The operational papers correctly keep a human in the loop, but they say less
about the workload created by false alarms. A method that produces one
nuisance flag per day on one series may produce hundreds on a large network.
At that scale, the cost is not only statistical; it is the attention of the
person reading the flags. Anscombe's old idea of treating rejection as an
insurance trade-off is still useful here \citep{anscombe1960}.

\subsection{Can a fault be separated from a real anomaly?}

Campbell and colleagues make this distinction explicit
\citep{campbell2013}. A sensor fault is something wrong with the instrument;
an anomaly may be a scientifically important event. Range checks are easy to
explain because impossible values are unlikely to be real. Spatial and model
checks are harder because a local storm or an unseen regime change can look
like a fault. This is a strong reason to keep flags and provenance instead
of automatically deleting data.

\subsection{Does the method transfer to a new site?}

Leigh and colleagues evaluate methods on expert-labelled data from a limited
number of settings \citep{leigh2019}. That is much stronger than testing on
unlabelled data, but a method that works at one site may still require new
thresholds or new training at another. The same issue appears in my
benchmark: the exact results are only valid for the simulated sampling rate,
noise and faults. Transferability should therefore be tested rather than
assumed.

These criticisms lead to one practical recommendation: QC should be treated
as a documented decision system, not merely as a list of functions. Each
check should state what evidence it uses, what assumptions it makes, where
its threshold came from, and what happens after it fires.


